% ======================================================================
% SECTION 8: AUTONOMOUS SUPPLY CHAIN AND IMP LOGISTICS
% File: sections/supply_chain.tex
% ======================================================================

The investigational medicinal product supply chain presents unique challenges for Physical AI
oncology trials: robotic pharmacy systems prepare and dispense doses, cold-chain logistics
span multiple robotic storage systems, and supply forecasting must account for both
pharmacological treatment schedules and robotic procedure timelines. The autonomous supply
chain system replaces the traditional CMC, supply, and logistics team with a dedicated supply
agent that manages drug and device release, inventory optimization, chain of custody tracking,
and integration with robotic pharmacy infrastructure.

\subsection{Supply Chain Agent}

The supply chain agent (supply\_agent) automates the complete IMP lifecycle from manufacturing
release through site-level dispensing and reconciliation. The agent manages four core
functions: lot release logic, inventory management, chain of custody, and shipment
optimization.

Lot release logic automates the review and approval of manufactured IMP lots for clinical
use. The agent evaluates certificates of analysis against predefined release specifications,
verifies stability data adequacy, confirms GMP compliance documentation, and generates
release decisions. Release decisions for novel modalities (cell therapies, gene therapies)
incorporate modality-specific release criteria including potency assays, sterility testing,
and viability requirements. Release decisions that fall outside validated parameters or
involve deviations from specification require human quality officer review.

Inventory management maintains real-time visibility across all supply chain nodes: central
warehouse, regional depots, site pharmacies, and robotic storage systems. The agent tracks
lot numbers, expiration dates, temperature history, and current location for every IMP unit.
Automated alerts trigger when inventory at any node falls below safety stock thresholds
or when expiration dates approach within the minimum remaining shelf-life window.

Chain of custody tracking provides an unbroken record of every IMP unit from manufacturing
through patient administration or destruction. Each custody transfer (warehouse to depot,
depot to site, site pharmacy to robotic system, robotic system to patient) generates an
electronic record with timestamp, location, handler identification, temperature reading,
and unit condition verification. This chain of custody satisfies 21 CFR Part 312
investigational drug accountability requirements and provides the audit trail needed for
regulatory inspection.

\subsection{CMC Lifecycle Management}

CMC lifecycle management spans the regulatory chemistry, manufacturing, and controls
requirements from pre-IND through commercial launch. The agent tracks the CMC lifecycle
through defined stages: pre-IND GMP-compliant pilot manufacturing, Phase I clinical supply,
process development and scale-up for Phase II/III, commercial manufacturing validation, and
post-approval lifecycle management per ICH Q12 \cite{ich-q12}.

GMP compliance verification ensures that manufacturing operations meet applicable GMP
standards: FDA 21 CFR Parts 210/211 for drug products, EU EudraLex Volume 4 for
investigational medicinal products manufactured for EU trials, and ICH Q7 for active
pharmaceutical ingredients. The agent tracks manufacturing site GMP certification status,
inspection history, and any regulatory actions that might affect supply continuity.

Manufacturing change management follows ICH Q12 principles for post-approval changes and
applies analogous principles to investigational product manufacturing changes. When
manufacturing process changes are proposed (scale changes, site transfers, raw material
source changes), the agent evaluates the change through a structured comparability
assessment framework. Changes classified as having potential to affect product quality
trigger IND amendments per \S312.30 with appropriate CMC updates.

\subsection{IRT-Driven Demand Forecasting and Distribution}

Interactive Response Technology drives the randomization and supply management integration
that determines drug supply requirements at each trial site. The supply agent interfaces
with the IRT system to obtain real-time enrollment data, treatment arm assignments, and
visit schedule projections that feed demand forecasting models.

Depot and cold-chain strategy optimization determines the number, location, and capacity
of regional distribution depots to minimize delivery time while controlling inventory
costs. For temperature-sensitive products (biologics, cell therapies), the agent optimizes
cold-chain logistics including packaging specifications, shipping route selection, and
temperature monitoring throughout the distribution chain. The optimization considers
seasonal temperature variations, transit time windows, and backup shipping options for
each route.

Temperature excursion detection and response automates the handling of temperature
deviations during storage and transit. When temperature monitoring data indicates an
excursion outside the validated range, the agent initiates an automated response sequence:
quarantine affected units, assess product impact based on excursion severity and duration,
request stability evaluation from the CMC team, and either release or destroy affected
units based on the assessment outcome. This automated response reduces the time between
excursion detection and disposition decision, minimizing supply disruption.

Site-level inventory prediction uses enrollment velocity data, treatment discontinuation
rates, and protocol visit schedules to forecast demand at each trial site. The agent
generates automated replenishment orders when predicted inventory falls below the site
safety stock level, accounting for lead times from depot to site. Prediction accuracy is
tracked as a KPI, with model parameters recalibrated monthly based on actual consumption
patterns.

\subsection{Robotic Supply Chain Integration}

Robotic supply chain integration extends the traditional pharmaceutical supply chain into
the Physical AI execution environment. Robotic pharmacy systems (Franka Emika collaborative
robots) perform dose preparation, verification, and dispensing under the autonomous supply
agent's coordination.

Table~\ref{tab:imp-supply} presents the IMP supply chain automation comparing traditional
and autonomous methods.

\begin{longtable}{L{2.0cm} L{2.5cm} L{2.8cm} L{1.8cm} L{2.5cm}}
\caption{IMP Supply Chain Automation}
\label{tab:imp-supply} \\
\toprule
\textbf{Process Step} & \textbf{Traditional Method} & \textbf{Autonomous Method} & \textbf{Agent Owner} & \textbf{Compliance Basis} \\
\midrule
\endfirsthead
\toprule
\textbf{Process Step} & \textbf{Traditional Method} & \textbf{Autonomous Method} & \textbf{Agent Owner} & \textbf{Compliance Basis} \\
\midrule
\endhead
\midrule
\multicolumn{5}{r}{\textit{Continued on next page}} \\
\endfoot
\bottomrule
\endlastfoot
Lot release & Manual review of CoA and stability data & Automated specification verification with deviation detection & supply\_agent & 21 CFR 211, ICH Q7 \\
Demand forecasting & Spreadsheet-based periodic forecasts & Real-time IRT-driven predictive models & supply\_agent & GCP, ICH E6(R3) \\
Depot distribution & Manual order management & Automated replenishment with cold-chain optimization & supply\_agent & GDP, EU Annex 15 \\
Site dispensing & Pharmacist manual preparation & Robotic preparation with automated verification & supply\_agent, robot\_execution\_gateway & 21 CFR 211, USP 797 \\
Temperature monitoring & Periodic manual checks & Continuous IoT sensor monitoring with automated alerts & supply\_agent & GDP, WHO TRS 961 \\
Chain of custody & Paper-based documentation & Electronic records with cryptographic provenance & supply\_agent, trust layer & 21 CFR Part 11, 21 CFR 312 \\
Reconciliation & Manual counting and documentation & Automated unit-level tracking and reconciliation & supply\_agent & ICH E6(R3), 21 CFR 312.62 \\
\end{longtable}

Table~\ref{tab:robotic-pharmacy} details the robotic pharmacy operations.

\begin{longtable}{L{2.2cm} L{2.2cm} L{2.5cm} L{2.5cm} L{2.2cm}}
\caption{Robotic Pharmacy Operations}
\label{tab:robotic-pharmacy} \\
\toprule
\textbf{Operation} & \textbf{Robot System} & \textbf{Safety Gate} & \textbf{Verification} & \textbf{Provenance Record} \\
\midrule
\endfirsthead
\toprule
\textbf{Operation} & \textbf{Robot System} & \textbf{Safety Gate} & \textbf{Verification} & \textbf{Provenance Record} \\
\midrule
\endhead
\midrule
\multicolumn{5}{r}{\textit{Continued on next page}} \\
\endfoot
\bottomrule
\endlastfoot
Dose calculation & Franka Emika cobot & Weight-based dosing verification against protocol & Independent algorithm cross-check & Patient ID, weight, dose, timestamp \\
Vial selection and preparation & Franka Emika cobot & Barcode verification of correct product, lot, and expiry & Visual confirmation system & Product ID, lot, expiry, vial serial \\
Reconstitution and dilution & Franka Emika cobot & Volume measurement verification within tolerance & Gravimetric verification & Volume, diluent, concentration, time \\
Final product labeling & Franka Emika cobot & Label content verification against order & OCR verification of printed label & Label image, verification result \\
Dispensing and handoff & Franka Emika cobot & Patient identity verification before release & Pharmacist or clinician sign-off & Recipient ID, handoff time, condition \\
\end{longtable}

The robotic pharmacy interface manages the handoff between the supply chain agent (which
determines what to prepare) and the robot execution gateway (which controls how the robot
performs the preparation). This separation of concerns ensures that supply decisions and
robotic execution are independently validated, providing defense-in-depth against
preparation errors. The Franka Emika collaborative robot platform achieves a Unification
Standard Level score of 88.75 in the national platform evaluation \cite{national-platform-paper},
reflecting strong performance in AI integration, cross-robot coordination, and clinical
trial collaboration dimensions.
